\documentclass{article}
\usepackage[margin=1.5cm]{geometry}
\usepackage{amsmath}
\usepackage{enumitem}
\usepackage{amssymb}

\setlist{nosep,leftmargin=*}

\begin{document}
\twocolumn

\title{COSC 120 Final Project Planning Worksheet}
\author{Prof. Jordan C. Hanson}
\date{}

\maketitle

\noindent
\textbf{The final project takes the place of the final exam.}
You may work \textbf{individually or with one partner}. Your goal is
to design, build, test, and present an extended Python program that
brings together the major ideas from the course.

\section{What Do You Want to Build?}

Describe a program that would be useful, interesting, scientific,
creative, or fun. It should solve a problem, model something, organize
information, automate a task, or create an interactive experience. \\ \vspace{0.25cm}

\textbf{Working title:}
\underline{\hspace{5.2cm}}

\textbf{Partner, if any:}
\underline{\hspace{4.7cm}}

\textbf{What will the program do?}

\vspace{0.75cm}

\textbf{Who would use it, and why?}

\vspace{0.65cm}

\section{Inputs, Outputs, and Data}

What information will enter the program, and what results will the program produce? \\ \vspace{2cm}

Which data structures seem useful?

\begin{itemize}
\item[$\square$] Lists
\item[$\square$] Dictionaries
\item[$\square$] Nested data structures
\item[$\square$] Objects created from classes
\end{itemize}

Give one example of data your program will store.

\vspace{0.55cm}

\section{Organize the Program}

Your project should make meaningful use of several ideas from the course, including functions, conditionals and loops, lists and/or dictionaries, and at least one class with useful attributes and methods.  Use modules and imports where appropriate.  Sketch a possible organization below: \\ \vspace{1cm}

\begin{verbatim}
main.py
    |
    +-- module1.py
    +-- module2.py
\end{verbatim}

What is an example of a class the program could use?  What would the class store, and what would it do?

\vspace{1cm}

\section{Think About an Algorithm}

An algorithm is a sequence of steps for solving a problem. Identify one part of your project where the program must search, sort, compare, simulate, calculate, update, or make a decision repeatedly. How could you organize the steps so the program does not perform unnecessary work?

\vspace{1cm}

Could a better data structure or algorithm make this part more efficient? Explain briefly.

\vspace{1cm}

\section{Define a Realistic First Version}

What is the \textbf{minimum working version} of your project?  That is, the simplest version that solves the problem?

\vspace{0.85cm}

If that version works, what feature would you add next?

\vspace{0.65cm}

\section{Plan the Presentation}

At the end of the semester, you will present your project to the
class. Your presentation should demonstrate the program and explain
how it works.  Plan to show:

\begin{itemize}
\item the problem or goal,
\item a live or recorded demonstration,
\item the important classes, functions, and data structures,
\item one interesting algorithm or design decision,
\item what worked, what was difficult, and what you learned.
\end{itemize}

\textbf{Before leaving today:} Be prepared to describe your project
idea in one or two sentences and identify the first piece of code you
would write.

\end{document}
